\chapter{Operations and Observability}
\label{ch:operations}

\section{Lifecycle selection and ordering}

\begin{lstlisting}[language=bash]
blackship up api
blackship up --all
blackship down production
blackship restart api
\end{lstlisting}

Selectors can be exact service names, full names, unambiguous prefixes, tags,
or \command{ALL}. Starting a target includes its dependencies; stopping uses the
configured cascade policy and reverse dependency order.

Preview broad actions before applying them:

\begin{lstlisting}[language=bash]
blackship restart --all --dry-run
blackship --format json restart --all --dry-run
\end{lstlisting}

\section{EVA live reconciliation}

EVA compares desired parameters with live values from \command{jls}, applies
updatable differences in one \command{jail\_set(JAIL\_UPDATE)} operation,
replaces managed rctl rules, and re-pins the cpuset. It runs pre/post EVA hooks
around the change.

\begin{lstlisting}[language=bash]
blackship eva api --dry-run
blackship eva api
\end{lstlisting}

Path, VNET mode, addresses, devfs ruleset, and other immutable identity or
isolation changes require a restart. EVA reports them but does not pretend to
apply them. A jail that is not running is skipped.

\section{Interactive and scripted access}

\begin{lstlisting}[language=bash]
blackship console api --user root
blackship exec api --user www -- id
blackship exec api --workdir /srv -e MODE=prod -- ./migrate
blackship exec api --clear-env -e HOME=/tmp -- env
blackship cp ./config.toml api:/usr/local/etc/app.toml
blackship cp api:/var/log/app.log ./app.log
\end{lstlisting}

Use \command{--} before the in-jail command so its flags are not parsed as
Blackship flags. Copy operations reject jail-to-jail transfers and path/symlink
escapes.

\section{Status, health, and statistics}

\begin{lstlisting}[language=bash]
blackship ps
blackship stats
blackship stats api
blackship health api
blackship health --watch --interval 5
blackship logs api --lines 200
blackship logs api --follow
\end{lstlisting}

Statistics aggregate per-process CPU and RSS from \command{jls} and
\command{ps} libxo JSON. Health status reflects configured checks and retry
state. \command{health} runs checks inside the jail and escalates to root.
The start period counts from the jail's recorded start, so a one-shot check on
a settled jail reports real results; within the period it reports
\setting{starting}, and a jail that is not running reports \setting{unknown}. Watch and follow modes are terminal streams; they are intentionally not
part of the single-envelope global JSON contract.

\section{Supervisor}

\command{blackship supervise} runs the warden loop. It monitors jail events,
executes health checks, applies recovery/backoff behavior, samples resource
usage, performs health-triggered QoS transitions, and can send a dead-man ping.
For boot operation, use the packaged rc script with the same manifest and data
directory as interactive administration.

The supervisor distinguishes intentional stops from failures and avoids
restarting a jail the operator deliberately took down. Each jail has an
independent one-for-one restart policy. Abnormal restarts back off from a
one-second base toward a 60-second ceiling, with jitter so simultaneous
failures do not restart in lockstep. The policy gives up when a jail exhausts
five restarts in a five-minute window rather than entering an endless failure
loop.

\section{QoS operations}

\begin{lstlisting}[language=bash]
blackship qos api --show
blackship qos api --steady
blackship qos api --startup
\end{lstlisting}

Only resources named by the target profile are swapped. Custom rules outside
the profile survive a phase transition. A one-shot \command{up} does not remain
resident to observe health; use the supervisor or switch manually.

\section{Audit history}

\begin{lstlisting}[language=bash]
blackship audit api
blackship --format json audit api
\end{lstlisting}

Audit JSONL is durable after the kernel jail, scope record, and even a reboot
are gone. The post-mortem summarizes event counts, lifetime, rollback steps,
termination reason, and sampled peak resources. Peaks are bounded by the sample
interval and are not exact continuous maxima.

\section{Orphan detection and garbage collection}

An out-of-band \command{jail -r}, crash, or hard reset can leave mounts,
interfaces, PF state, leases, rctl rules, or datasets behind. Garbage collection
compares durable scopes with live host state:

\begin{lstlisting}[language=bash]
blackship gc --dry-run
blackship gc
blackship gc --force
\end{lstlisting}

The normal mode reclaims only resources Blackship can prove it owns. Datasets
not created by Blackship and non-ephemeral jails are reported rather than
destroyed. \command{--force} expands authority and demands careful review.

\section{Telemetry}

When configured, lifecycle events and racct gauges are exported over OTLP/HTTP.
Telemetry runs out of band: collector failure emits a warning but cannot make a
jail lifecycle operation fail. This prevents observability from becoming a
control-plane dependency.

\section{Operational checklist}

Before a change window:

\begin{enumerate}
  \item run \command{blackship check};
  \item inspect \command{ps}, \command{health}, and recent \command{audit};
  \item snapshot stateful ZFS roots or take a portable export;
  \item generate the relevant dry-run plan;
  \item verify privilege mode and selected targets; and
  \item ensure an operator can access the host without depending on the jail
    being changed.
\end{enumerate}

After the change, verify health, port forwards, dependency status, and audit
termination/start records. Run a non-forced GC dry run to catch residual state.
